\section{A finite positional attention realization}\label{sec:implementation-proof}

This section realizes the approximate-statistic interface by an explicitly
specified finite causal attention/Boolean--ReLU program. Its purpose is to
verify the observation cut and its costs. A sequential counter on the same
input is a stronger computational baseline, described in
Section~\ref{sec:impl-comparisons}. Uniform aggregation of indicators and
positional copying are established operations
\citep{WeissGoldbergYahav2021,EdelmanEtAl2024}. Moreover, the earlier
finite-bank construction \citep{BCJ001} already contains a paid uniform
positional bank, a real-error half-gate analysis, and a finite sequential
evaluator. The construction here specifies singleton-copy heads, direct
$J/C/D$ rails, nearest rounding, a canonical post-rounding replacement cut,
and an explicit Boolean--ReLU suffix with a separately stated hybrid shape.
It does not establish a new statistical precision phenomenon or a learned
mechanism.

\subsection{The finite model and its guarantee}\label{sec:impl-statement}

For this section only, let $M\ge1$ and $h,b\ge0$ be integers. The parameter
$b$ denotes \emph{output} precision; $r$ remains the raw-window width used
elsewhere. Set
\begin{equation}\label{eq:impl-scales}
 n=M+3,\qquad q=M+2,\qquad Q=2^h,\qquad
 \ell_a=\left\lceil\log_2(n+1)\right\rceil,\qquad
 p=\ell_a+b+h+32.
\end{equation}
The address width $\ell_a$ is distinct from the logarithmic competitive
factor $\ell$ in Theorem~\ref{thm:finite-window}. The input consists of the
$M+2$ binary symbols $Z_0,\ldots,Z_{M+1}$ and one synthetic zero token at
position $q$. No one-hot representation, endpoint label, local pattern
feature, or count is supplied without construction.

For $u=(a,c)\in\{0,1\}^2$, define on an arbitrary binary word
\begin{align}
 J'_u&=\sum_{i=0}^{M-1}\ind\{(Z_i,Z_{i+1})=u\},\notag\\
 C'_u&=\sum_{i=0}^{M-1}\ind\{(Z_i,Z_{i+1})=u,\ Z_{i+2}=1-a\},\label{eq:impl-counts}\\
 D'_u&=\sum_{i=0}^{\lfloor M/2\rfloor-1}
                \ind\{(Z_i,Z_{i+1})=u\}.\notag
\end{align}
An empty sum is zero. A \emph{finite attention head} in the present model
has explicit public structural masks and exact bounded-word evaluation of
its allowed-key sum, followed by a specified rounding operation. The
integer work of that evaluator is charged separately from the number of
attention layers. Pointwise maps consist of affine maps and ReLUs, with
exact finitely represented coefficients and explicit live-wire holds. This
is a hybrid computational model, not unit-cost real arithmetic or a
finite-precision hardware softmax claim.

\begin{theorem}[Finite realization of the positional bank]\label{thm:finite-realization}
For each $M,h,b$ above there is a fixed causal finite program, depending
only on these public parameters, with the following properties.
\begin{enumerate}
\item It has two attention layers and three heads in total, with at most
 two heads in one layer, query/key dimension one, maximum value dimension
 twelve, and score cap zero. Shared pointwise Boolean--ReLU maps have
 depth $O(p^2)$ and width $O(p)$ per token. The integer attention-evaluator
 work is additional and is charged below.
\item On every binary input word, it acquires the exact pair
 $(Z_M,Z_{M+1})$ and produces twelve master numerators
 $B_{u,A}\in\{0,\ldots,Q\}$, $A\in\{J,C,D\}$, satisfying
 \begin{equation}\label{eq:impl-bank-error}
 \left|\frac{B_{u,J}}Q-\frac{J'_u}M\right|,
 \quad\left|\frac{B_{u,C}}Q-\frac{C'_u}M\right|,
 \quad\left|\frac{B_{u,D}}Q-\frac{D'_u}M\right|
 \ \le\ 2^{-h-1}.
 \end{equation}
 The inequalities are simultaneous and deterministic.
\item At the specified cut, the only data-dependent downstream inputs
 are those twelve numerators and the exact endpoint pair. Replace all
 twelve masters jointly by arbitrary canonical integers
 $B^*_{u,A}\in\{0,\ldots,Q\}$ satisfying
 \begin{equation}\label{eq:impl-port}
 \max_{u,A}\frac{|B^*_{u,A}-B_{u,A}|}{Q}\le\delta.
 \end{equation}
 The resulting normalized coordinates have error at most
 $\eta=2^{-h-1}+\delta$ from \eqref{eq:impl-counts} divided by $M$.
 Replacements may be correlated and may depend on the observed word.
 Endpoint bits, metadata, comparators, and subsequent arithmetic remain
 exact; this is the entire numeric fault port.
\item The readout selects the actual endpoint row, clips its $C,D$
 coordinates to $[0,J]$, and applies the gate $J>0$, $512D\ge J$.
 Its output is a binary probability distribution with denominator $2^b$,
 within $2^{-b}$ total variation of the corresponding exact rational gate
 output, including empty rows, gate equality, $h=0$, and $b=0$.
\item In the experiment of Theorem~\ref{thm:approximate-gate}, with
 $M\ge64$ and supplied certificate $1/M$, its actual-terminal risk obeys
 \begin{equation}\label{eq:impl-risk}
 R\le\min\left\{1,\frac{131072}{\sqrt M}
       +6433056\frac{k}{M}
       +1072176\bigl(2^{-h-1}+\delta\bigr)+2^{-b}\right\}.
 \end{equation}
 The program is not supplied $k$ or $\delta$.
\end{enumerate}
A complete sparse-list/sparse-circuit hybrid realization has conservative
online bitwork $O(np^4)$ and total online peak storage
$O(p^4+n\ell_a+np^2)$ bits. A complete dense-mask/dense-circuit alternative
has online bitwork $O(n^2p^2+np^6)$ and total peak storage
$O(p^5+n^2+n\ell_a+np^2)$ bits. These are alternative full
representations, not interchangeable entries in one cost calculation.
Section~\ref{sec:impl-ledgers} gives construction, input, addressing,
mask, hold, cleanup, and output costs and a separately charged all-ReLU
unrolling.
\end{theorem}

The endpoint in this theorem is computed from the input, not appended as
an additional observation. The fault restriction in
\eqref{eq:impl-port} is substantive: off-grid analog errors, corrupted
endpoint bits, accumulator errors, and comparator faults are not covered.
The construction and the following proof do not assume that the rounded
or replaced bank is the count table of any binary word.

\subsection{Causal masks and the twelve direct rails}\label{sec:impl-prefix}

Write $\rho(x)=\max\{x,0\}$. Initially
\[
 \operatorname{raw}_t=Z_t\quad(0\le t\le M+1),\qquad
 \operatorname{raw}_q=0.
\]
Public position information provides the query-type flag
$U_t=\ind\{t=q\}$ and the half flag
\begin{equation}\label{eq:impl-half-flag}
 F_t=\ind\{2\le t<2+\lfloor M/2\rfloor\}.
\end{equation}
These flags are generated by public integer comparisons. Positions and
flags cost $O(n\ell_a)$ stored bits and their generation is charged below.

\paragraph{First layer: two singleton copies.}
Both heads have $Q_t=K_t=0$ and the one-dimensional value
$\operatorname{raw}_t$. Their permitted-key sets are
\begin{equation}\label{eq:impl-copy-masks}
 \mathcal M^{(1)}_t=\begin{cases}\{t-1\},&t\ge1,\\\{t\},&t=0,
                       \end{cases}
 \qquad
 \mathcal M^{(2)}_t=\begin{cases}\{t-2\},&t\ge2,\\\{t\},&t<2.
                       \end{cases}
\end{equation}
Each row is nonempty and causal. Its only normalized attention weight is
exactly one, so the returned bits, denoted $l_{1,t}$ and $l_{2,t}$,
are exact copies. Raw bits and both copied bits are retained through paid
identity holds. In particular,
\begin{equation}\label{eq:impl-copy-identities}
 (l_{2,i+2},l_{1,i+2},\operatorname{raw}_{i+2})
       =(Z_i,Z_{i+1},Z_{i+2})\quad(0\le i<M),\qquad
 (l_{2,q},l_{1,q})=(Z_M,Z_{M+1}).
\end{equation}
The fallback copies at positions zero and one never enter the query's
statistic. These masks are exact structural exclusions. They are not
approximations obtained by assigning large scores in a dense causal row.

\paragraph{Local pointwise features.}
For each $u=(a,c)$ in lexicographic order, form
\begin{align}
 R_{t,u}&=\bigl(1-(l_{2,t}\mathbin{\mathrm{XOR}}a)\bigr)
              \mathbin{\mathrm{AND}}
              \bigl(1-(l_{1,t}\mathbin{\mathrm{XOR}}c)\bigr),\notag\\
 V_{t,u,J}&=R_{t,u},\qquad
 V_{t,u,C}=R_{t,u}\mathbin{\mathrm{AND}}
                    (\operatorname{raw}_t\mathbin{\mathrm{XOR}}a),
 \qquad
 V_{t,u,D}=R_{t,u}\mathbin{\mathrm{AND}}F_t.
 \label{eq:impl-features}
\end{align}
These are twelve binary rails, not an eight-successor histogram with an
additional reconstructed row sum. On the legal key $t=i+2$, substitution
in \eqref{eq:impl-copy-identities} gives respectively
\begin{equation}\label{eq:impl-feature-identities}
 \ind\{(Z_i,Z_{i+1})=u\},\qquad
 \ind\{(Z_i,Z_{i+1})=u,\ Z_{i+2}=1-a\},\qquad
 \ind\{(Z_i,Z_{i+1})=u,\ i<\lfloor M/2\rfloor\}.
\end{equation}
Thus the $C$ rail uses precisely the designated successor. Each feature
uses only its own token's local copies and the public half flag. In
particular, no endpoint-conditioned query is broadcast to historical
tokens.

\paragraph{Second layer: one uniform head.}
This head again has one-dimensional $Q=K=0$, and its value is the
twelve-dimensional vector in \eqref{eq:impl-features}. Its mask is
\begin{equation}\label{eq:impl-uniform-mask}
 \mathcal M^{(3)}_q=\{2,3,\ldots,M+1\},\qquad
 \mathcal M^{(3)}_t=\{t\}\quad(t\ne q).
\end{equation}
The query has exactly $M$ allowed keys. The key $M+1$ includes the last
legal departure triple, whereas the synthetic query $q$ is excluded.
The terminal pair is not an additional departure. Every other row uses a
causal singleton self mask; all those rows are evaluated and their
outputs are subsequently cleared. This choice is part of the construction
and is not a dense causal-prefix restriction.

All allowed scores equal zero, so $\exp(0)=1$ and the ideal normalized
head is a uniform average. There is no approximation of a logarithm,
exponential, anchor inversion, or content match. For each row with $L$
allowed keys, sum each binary rail into an \emph{exact} integer
$S\in\{0,\ldots,L\}$ and return its nearest grid numerator
\begin{equation}\label{eq:impl-round}
 \mathcal R_{h,L}(S)
 =\left\lfloor\frac{2QS+L}{2L}\right\rfloor.
\end{equation}
The denominator $L$ is public, obtained from the enumerated mask, and is
never zero. The floor in this formula will be implemented by the finite
division circuit below; it is not a primitive real-number operation.
At $q$, summing \eqref{eq:impl-feature-identities} proves the all-word
identity
\begin{equation}\label{eq:impl-exact-sums}
 (S_{q,u,J},S_{q,u,C},S_{q,u,D})=(J'_u,C'_u,D'_u).
\end{equation}
Equation~\eqref{eq:impl-round} rounds $QS/L$ to its nearest integer, with
ties upward. It maps $0$ to $0$ and $L$ to $Q$, and
\[
 \left|\frac{\mathcal R_{h,L}(S)}Q-\frac SL\right|
 \le\frac1{2Q}.
\]
Together with \eqref{eq:impl-exact-sums}, this proves
\eqref{eq:impl-bank-error} for every binary word, including $M=1$ and
$h=0$. Counts are accumulated before this single rounding. Repeated
$h$-grid rounding during accumulation would be a different algorithm.

\subsection{Exact Boolean and ReLU arithmetic}\label{sec:impl-arithmetic}

All words use little-endian bit rails, with enough leading zeros for the
stated operation. For bits $x,y$, the identities
\begin{align}
 x\mathbin{\mathrm{AND}}y&=\rho(x+y-1),&
 \mathrm{NOT}(x)&=\rho(1-x),\notag\\
 x\mathbin{\mathrm{XOR}}y&=\rho(x-y)+\rho(y-x),&
 x\mathbin{\mathrm{OR}}y&=x+y-\rho(x+y-1)
 \label{eq:impl-boolean}
\end{align}
are exact. Each expression is a constant-size affine/ReLU circuit. A
coordinatewise multiplexer is
\begin{equation}\label{eq:impl-mux}
 \operatorname{mux}(s,x,y)
 = (s\mathbin{\mathrm{AND}}x)
      \mathbin{\mathrm{OR}}
   ((\mathrm{NOT}(s))\mathbin{\mathrm{AND}}y).
\end{equation}
It equals $x$ when $s=1$ and $y$ when $s=0$. Thus no data-dependent
selection below is free dynamic addressing.

\paragraph{Addition, subtraction, and comparison.}
For input bits $x_i,y_i$ and carry $c_i$, a full adder has
\[
 v_i=x_i\mathbin{\mathrm{XOR}}y_i,\qquad
 s_i=v_i\mathbin{\mathrm{XOR}}c_i,\qquad
 c_{i+1}=(x_i\mathbin{\mathrm{AND}}y_i)
             \mathbin{\mathrm{OR}}
            (v_i\mathbin{\mathrm{AND}}c_i).
\]
Induction on $i$ gives ordinary binary addition. A ripple chain on
$w$-bit words uses $O(w)$ depth and active gates, with $O(w)$ live rails.
Retain the final carry separately. Subtract $y$ from $x$ by adding
$x+\mathrm{NOT}(y)+1$. Since this integer is $x-y+2^w$, its final carry
is one exactly when $x\ge y$, and its low $w$ bits are $x-y$ modulo
$2^w$. Hence the same chain implements unsigned comparison and
subtraction. Public shifts are explicit fixed wires with zero padding;
their storage and holds are charged.

\begin{lemma}[Finite restoring division]\label{lem:impl-division}
For a positive $w$-bit denominator and a $w$-bit dividend, exact quotient
and remainder can be computed with $w$ fixed public iterations, $O(w^2)$
active Boolean/ReLU gates and depth, and $O(w)$ live rails. A layered
representation that holds all live rails at every step has at most
$O(w^3)$ gates and identity holds.
\end{lemma}
\begin{proof}
Initialize remainder $R=0$. Scan the dividend's bits from most significant
to least. Given the next bit $x_i$, first form $\widetilde R=2R+x_i$.
Subtract the denominator $d$ using the preceding circuit. Its final carry
$c$ says whether $\widetilde R\ge d$. Emit $c$ as the next quotient bit,
and replace the remainder by
$\operatorname{mux}(c,\widetilde R-d,\widetilde R)$.

After scanning any prefix $P$ of the dividend, the invariant is
$P=Q_Pd+R$ with $0\le R<d$, where $Q_P$ is the emitted quotient prefix.
For the next bit, $2R+x_i<2d$, so one subtraction suffices, proving the
invariant inductively. Also $R\le P$, and therefore
$\widetilde R\le2P+x_i$, the next scanned dividend prefix. That prefix
fits in the original $w$ bits. Thus the shifted remainder fits without an
uncharged extra magnitude bit. The subtraction's carry is one explicit
additional rail. At the final iteration the invariant is the desired
Euclidean division identity. Each of $w$ iterations costs $O(w)$ active
gates and depth. Holding $O(w)$ live rails through $O(w^2)$ layers gives
the stated conservative $O(w^3)$ bound.
\end{proof}

The same sequence is a total Boolean circuit even on denominator zero;
its arbitrary zero-denominator result will never be used. The readout
below instead selects a strictly positive denominator before invoking
this circuit.

\paragraph{Word bounds for the attention evaluator.}
A feature sum needs $\lceil\log_2(M+1)\rceil$ bits; $h$ bits alone do not
suffice. The rounding dividend and divisor satisfy
\begin{equation}\label{eq:impl-prefix-wordbounds}
 2QS+L\le 2^{h+1}M+M<2^{h+\ell_a+2},\qquad
 2L\le2M<2^{\ell_a+1}.
\end{equation}
The same bounds cover singleton rows because $M\ge1$. Use
$w_{\rm att}=\ell_a+h+16$ bits, together with explicitly retained carries.
The accumulation, the $(h+1)$-place shift, addition of $L$, and restoring
division all fit. In particular, the accumulation is exact, and the
rounding numerator is computed without overflow. This implements
\eqref{eq:impl-round} entirely by the preceding finite gadgets.

\subsection{Packed and scalar versions of the bank}\label{sec:impl-extraction}

The bank has two possible physical representations, with the same declared
observation. In a packed version, each master is stored as its $h+1$
numerator bits. Reading these existing bits and padding them with zeros
is paid fixed wiring. The high bit is necessary to represent $B=Q$, that
is, the normalized value one.

In a scalar version, the head returns one exact grid scalar
$H=B/Q$ per rail. The head's conversion from numerator bits to this
scalar is the paid dyadic sum
\[
 H=\sum_{j=0}^{h}2^{j-h}B_j.
\]
It uses $O(h+1)$ terms with $O(h+1)$-bit coefficients. A scalar is
\emph{not} treated as exposing its bits at no cost.

\begin{lemma}[Paid extraction from a canonical scalar grid]
\label{lem:impl-extraction}
For $H\in\{0,1/Q,\ldots,1\}$, the numerator bits of $QH$ are extracted
exactly by a ReLU circuit of depth and active-gate count $O(h+1)$, with
$O((h+1)^2)$ conservative holds and $O(h+1)$-bit coefficients and scalar
words. No floor on arbitrary real inputs is used.
\end{lemma}
\begin{proof}
Initialize $A=QH$ and, in the fixed order $j=h,h-1,\ldots,0$, set
\begin{equation}\label{eq:impl-extraction}
 B_j=\rho\bigl(A-(2^j-1)\bigr)-\rho(A-2^j),\qquad
 A\leftarrow A-2^jB_j.
\end{equation}
Initially $A$ is an integer in $[0,Q]$. At the beginning of step $j$ it
is an integer in $[0,2^{j+1}-1]$. If $A<2^j$, both ReLUs in
\eqref{eq:impl-extraction} vanish. If $A\ge2^j$, their difference is
exactly one. The update leaves an integer in $[0,2^j-1]$. Induction gives
all the bits and final remainder zero. At $H=1$, the top bit $B_h$ is
one and all remaining bits are zero. At $h=0$, the single step extracts
either zero or one. Each step has constant active cost; retaining all
emitted bits through later steps costs $O((h+1)^2)$ identity holds.
Thresholds, scaling coefficients, and exact scalar values have
$O(h+1)$-bit descriptions.
\end{proof}

For the scalar realization the replacement supplies a canonical scalar
$H^*=B^*/Q$, and extraction occurs \emph{after} replacement. Consequently
it reveals only the replaced master, not its earlier value. On off-grid
inputs the recurrence need not produce bits; such inputs are outside the
fault model. Twelve parallel extractions, their scaling, their scalar
head decodes, and their holds fit the $p$-based resource bounds below.
Both representations therefore reach the same padded numerator rails
by paid operations.

\subsection{Endpoint selection and a total finite readout}
\label{sec:impl-readout}

Let $(e_0,e_1)=(Z_M,Z_{M+1})$ be the two copied endpoint bits. For each
$u=(a,c)$ form the selector
\[
 s_u=\bigl(1-(e_0\mathbin{\mathrm{XOR}}a)\bigr)
          \mathbin{\mathrm{AND}}
       \bigl(1-(e_1\mathbin{\mathrm{XOR}}c)\bigr).
\]
Exactly one selector equals one. For each of the three coordinate types,
start a zero word and traverse the four rows in fixed order, replacing
it by $\operatorname{mux}(s_u,B^*_{u,A},\text{current word})$.
The resulting words are the selected canonical numerators. This is a
constant number of word muxes; it does not use an uncharged lookup indexed
by the endpoint.

Denote the selected words by $J,C_0,D_0$, and set
\begin{equation}\label{eq:impl-clip-gate}
 C=\min\{C_0,J\},\qquad D=\min\{D_0,J\},\qquad
 G=\ind\{J\ge1\}\mathbin{\mathrm{AND}}\ind\{512D\ge J\}.
\end{equation}
Each minimum is implemented by comparison and a mux. Lower clipping is
automatic from unsigned canonical inputs. Multiply $D$ by $512$ with a
nine-place shift and compare exactly. Gate equality therefore accepts,
whereas $J=0$ rejects, irrespective of the second comparison.

Choose the safe denominator and quotient by
\begin{equation}\label{eq:impl-output}
 J_{\rm safe}=\operatorname{mux}(G,J,1),\qquad
 v=\left\lfloor\frac{2^bC}{J_{\rm safe}}\right\rfloor,\qquad
 m=\operatorname{mux}(G,v,0).
\end{equation}
The divisor is always positive. The shift and restoring division implement
the displayed floor. Even on a rejected row the pre-mux quotient is at
most $2^{h+b}$. On acceptance, $0\le C\le J$ gives $0\le m\le2^b$.
On rejection $m=0$. Compute the complement numerator $2^b-m$ by exact
subtraction and route the designated probability using the endpoint bit:
\begin{equation}\label{eq:impl-symbol-route}
 (P_0,P_1)=
 \begin{cases}
 (2^b-m,m),&e_0=0,\\
 (m,2^b-m),&e_0=1.
 \end{cases}
\end{equation}
Again this is a pair of bitwise muxes. The actual output is
$(P_0/2^b,P_1/2^b)$, always in the binary simplex. A dyadic decoder
$\sum_{i=0}^b 2^{i-b}(P_a)_i$ uses paid coefficients of $O(b+1)$ bits.

On acceptance, flooring gives
$0\le C/J-m/2^b<2^{-b}$; on rejection the rational and finite outputs
both put zero on the designated symbol. Binary total variation is the
absolute difference of that probability. This proves the output error
bound, including $b=0$, without a gate-margin assumption.

All readout arithmetic can use $w_{\rm out}=h+b+16$ bits. Indeed,
$512D\le2^{h+9}$, $2^bC\le2^{h+b}$, the pre-mux quotient is no larger
than this dividend, and the final probability numerators are at most
$2^b$. Constants, intermediate carries, and padded words are bounded by
$2p$ rails. Together with \eqref{eq:impl-prefix-wordbounds}, these bounds
prove the absence of overflow throughout the construction.

The pointwise feature and suffix maps are shared and executed at
\emph{all} $n$ tokens. The public flag $U_t$ clears packed bank and endpoint
inputs at nonquery tokens using muxes. For a scalar bank rail
$H_t\in[0,1]$, the paid gate $\rho(H_t+U_t-1)$ retains $H_t$ at the
query and is zero elsewhere. Thus nonquery banks are cleared before the
cut without extracting the query's bits early. Scalar extraction after
replacement is executed at all tokens, with zero inputs away from the
query. The shared suffix still computes on those zero inputs; its work
is not charged only at $q$. Only the query's output is retained.

\subsection{Isolation of the replacement cut and statistical composition}
\label{sec:impl-cut}

The master cut is after the twelve query coordinates have been rounded
and before endpoint-conditioned selection, comparison, division, or any
readout fanout. Every downstream bank copy originates from the replaced
master. In the scalar realization it also precedes the extraction in
Lemma~\ref{lem:impl-extraction}. The endpoint bits and public
$h,b$, type flags, and layout remain exact. Nonquery banks and endpoint
rails are cleared by the public query flag.

Before the cut is made available to the suffix, raw-input arrays, lag-copy
arrays, local-feature arrays, attention accumulators, rounding dividends
and divisors, intermediate quotients and remainders, and any clean-master
duplicates are released or overwritten. The suffix has no later attention
layer and no raw-memory read. Registers retaining $M$ or layout contain
only public, data-independent values; they do not encode any discarded
count. An equivalent interface signature is simply
\[
 \bigl((B^*_{u,A})_{u,A},\ (e_0,e_1),\ h,b\bigr)
       \longmapsto (P_0,P_1).
\]
There is no original-bank handle in this signature. All initialization,
clearing, and identity holds are included in the resource accounts.

This is a logical access restriction on the program, not a claim of
secure physical erasure against a hostile host retaining its own copies.
The external fault injector is not counted as free neural hardware: its
construction and runtime are outside the nominal program. The theorem
charges the program on either side of the stated replacement port, not
an implementation of the adversary.

\begin{proof}[Completion of the proof of Theorem~\ref{thm:finite-realization}]
The copied-bit identity \eqref{eq:impl-copy-identities}, the feature
identity \eqref{eq:impl-feature-identities}, and the exact mask
\eqref{eq:impl-uniform-mask} prove the twelve all-word sums
\eqref{eq:impl-exact-sums}. Exact finite division implements their nearest
rounding, proving \eqref{eq:impl-bank-error}. The replacement triangle
inequality yields
\[
 \left|\frac{B^*_{u,A}}Q-\frac{A'_u}M\right|
 \le\frac1{2Q}+\delta=\eta
 \qquad(A\in\{J,C,D\}).
\]
For analysis only, multiply the normalized replacement coordinates by
$M$, obtaining approximate counts with simultaneous error at most
$e=M\eta$. No such multiplication by $M$ is needed or supplied to the
actual readout. The canonical $J$ already lies in $[0,Q]$, and
\eqref{eq:impl-clip-gate}, after the common positive scaling $M/Q$,
is exactly the sequential clipping and gate of
Theorem~\ref{thm:approximate-gate}. Both the ratio and the comparison are
unchanged by this scaling. Inconsistent replacement coordinates are
permitted because that theorem does not require count realizability.

The protected endpoint equals the actual terminal pair and has been
acquired by the first layer. Equations~\eqref{eq:impl-output}--
\eqref{eq:impl-symbol-route} give output error at most $\xi=2^{-b}$.
Applying Theorem~\ref{thm:approximate-gate} with $e=M\eta$ and that $\xi$
proves \eqref{eq:impl-risk}. This is the only stochastic step. The
all-word interface identities need neither a source law nor an attack
independence assumption. The statistical composition retains the source,
strict-edit, and observation-measurability assumptions of the approximate
gate theorem. The finite circuits above give the claimed shape; the next
subsection verifies their full cost bounds.
\end{proof}

In particular, if $h,b\ge\lceil\tfrac12\log_2M\rceil$ and
$\delta=O(M^{-1/2})$, \eqref{eq:impl-risk} is
$O(M^{-1/2}+k/M)$. Theorem~\ref{thm:bank-precision} calibrates the clean,
fault-free \emph{fixed bank} obtained at this cut. It does not make the
present word widths or total program precision necessary. The phrase
``half-log precision'' can refer only to the fractional precision $h$ of
these normalized coordinates, not to all arithmetic, weights, addresses,
input storage, or total memory.

\subsection{Complete resource accounts}\label{sec:impl-ledgers}

All following bounds are constructive asymptotic upper bounds from the
gate recipes and fixed public schedules, with absolute constants. They
are not timing measurements or assertions that a conventional softmax
kernel executes the integer evaluator at unit cost. We count a bit
operation at the elementary Boolean/arithmetic level and charge the
selected program representation, addresses, and all holds. Only one
representation is selected at a time.

\paragraph{Hybrid shape and coefficient descriptions.}
There are two attention layers: two one-bit copying heads followed by one
twelve-value uniform head. All scores are zero; query/key dimensions are
one. A constant number of $O(p)$-word selections, comparisons, and
restoring divisions gives shared pointwise depth $O(p^2)$ and width
$O(p)$. The twelve scalar extractions, when present, add only $O(h+1)$
depth and fit that width. Attention's exact integer evaluator is charged
as work below and is not hidden in this small hybrid depth.

Individual coefficients use $O(p)$ bits. Boolean coefficients are bounded
constants, scalar-extraction coefficients use $O(h+1)$ bits, and output
dyadics use $O(b+1)$ bits. The shared pointwise circuit has $O(p^2)$ active
gates. Explicitly holding $O(p)$ live rails through all layers gives
$O(p^3)$ sparse gate/hold records. Their addresses and coefficients cost
at most $O(p)$ bits each, hence $O(p^4)$ stored sparse parameter bits.
The dense alternative has $O(p^2)$ layers, $O(p)$ width, and therefore
$O(p^4)$ coefficient slots including zeros and biases. Storing all
coefficients costs $O(p^5)$ bits. Constantly many head-projection matrices
fit these bounds. Finite evaluator code and bounded-word templates cost
at most $O(p^4)$ bits and are included, not left as an external oracle.

\paragraph{Masks, public positions, and input.}
Binary input acquisition costs $M+2$ bits. The synthetic zero is public.
For positions, query type, half flags, row boundaries, and masks, the two
complete choices are:
\begin{enumerate}
\item \emph{Active-key lists.} The first two heads each contain $n$
 stored keys. The third contains $n-1+M$ keys. Thus the exact number of
 stored key occurrences is $3n-1+M$. Each key and row boundary uses
 $O(\ell_a)$ bits, giving $O(n\ell_a)$ mask/list storage, including
 positions and flags. Enumerating indices, comparisons, routes, and row
 boundaries costs at most $O(n\ell_a^2)$ bitwork. Each stored route is
 read and charged during execution. This is explicitly paid sparse
 structural masking, not free addressing.
\item \emph{Literal dense masks.} Three $n$-by-$n$ masks cost exactly
 $3n^2$ mask bits, in addition to $O(n\ell_a)$ metadata. Generating and
 writing them costs at most $O(n^2\ell_a)$ bitwork. A dense online scan
 costs at most $O(n^2p)$ even though many permitted-key sets are
 singletons. Converting active lists to this representation pays these
 dense generation and storage costs.
\end{enumerate}

\paragraph{Finite attention arithmetic.}
There are $O(n)$ allowed key occurrences in the active-list version.
Each copied value is a bit. For the uniform head, each of twelve rails
per key uses an $O(p)$-bit addition, and every token performs twelve
$O(p^2)$ rounding divisions. Thus the fixed-mask finite attention
evaluation costs $O(np^2)$ elementary bitwork, including copies,
projections, and clearing. No logarithm, exponential, or source-specific
cache is generated. Dense-mask scanning adds the cost just stated.
Controllers and word addresses occupy $O(p)$ bits; the larger complete
runtime bounds below also allow a full $O(p)$ overhead per addressed
record lookup.

\paragraph{Two complete online implementations.}
Sparse execution of all pointwise maps, including every live hold, costs
$O(np^3)$ elementary bitwork. Combining this with active-list attention,
then charging up to $O(p)$ per sparse record access, gives the conservative
unconditional sparse/sparse online bound
\begin{equation}\label{eq:impl-sparse-work}
 T_{\rm sparse}=O(np^4).
\end{equation}
Its stored parameters and code are $O(p^4)$ bits, its masks and public
metadata are $O(n\ell_a)$ bits, and its observed input costs $M+2$ bits.
Dense coefficient counts are not substituted into this runtime.

For the dense-mask/dense-circuit alternative, coefficient-bit operations
are charged even for zero terms. Dense QKV projections, all-token shared
maps, mask scans, and word operations give the conservative bound
\begin{equation}\label{eq:impl-dense-work}
 T_{\rm dense}=O(n^2p^2+np^6).
\end{equation}
Its parameters and code cost $O(p^5)$ bits, with masks and metadata
$O(n^2+n\ell_a)$ bits. The quadratic term is a safe allowance for this
dense implementation, not a lower bound or an assertion that the sparse
implementation performs quadratic work.

\paragraph{Activations, cleanup, and total online peak.}
With bit rails and packed words, the working activation arrays require
$O(np)$ bits. A conservative scalar-channel implementation allocates
$O(p)$ physical bits to each of $O(p)$ live channels per token, hence
$O(np^2)$ bits. These allowances include Q/K/V arrays, feature residuals,
lag holds, twelve pre-clearing bank words per token, endpoint holds,
controller, and arithmetic scratch. Input costs $O(n)$ additional bits;
output costs $O(b+1)$ packed bits, or $O((b+1)p)$ in the conservative
scalar allocation. The explicit +1 covers $b=0$ and the numerator $2^b$.
Initializing, overwriting, and releasing these arrays is charged in the
online bounds. Only the latest live layer is retained: a training tape or
all-layer retention would be an additional cost.

Total online peak is the sum of simultaneously resident program and
parameters, masks and metadata, observed input, QKV and residual arrays,
banks, controller and scratch, and output. With conservative scalar
allocation it is bounded by
\begin{align}
 S_{\rm sparse}&=O(p^4+n\ell_a+np^2),
                  \label{eq:impl-sparse-memory}\\
 S_{\rm dense}&=O(p^5+n^2+n\ell_a+np^2).
                  \label{eq:impl-dense-memory}
\end{align}
The $12(h+1)$ bank bits and two endpoint bits are only the cut payload,
not these total-memory quantities. At $h$ of order $\tfrac12\log_2M$,
exact accumulators, addresses, and public horizon words still require
order $\log M$ bits, and the resident input and paid mask storage remain.

\paragraph{Offline construction and peak storage.}
Generate the ripple circuits, fixed division schedules, muxes, scalar
extractions, and all identity holds in $O(p^5)$ conservative bitwork with
$O(p^4)$ sparse scratch/record storage, and write the $O(p^4)$ sparse
program bits. Sorting, if used, involves $O(p^3)$ records of $O(p)$ bits
and $O(\log p)$ comparison overhead, which fits the same work allowance.
The dense representation additionally writes $O(p^5)$ coefficient bits,
including zero slots, within $O(p^5)$ bitwork. Add
$O(n\ell_a^2)$ work for sparse lists or $O(n^2\ell_a)$ work for dense
masks and their public metadata. Accordingly, safe full construction
work bounds are
\begin{equation}\label{eq:impl-offline-work}
 O(p^5+n\ell_a^2)\quad\text{(sparse)},\qquad
 O(p^5+n^2\ell_a)\quad\text{(dense)}.
\end{equation}
Offline peak includes partially generated stored program and masks plus
live generation scratch. It is bounded by the corresponding total in
\eqref{eq:impl-sparse-memory} or \eqref{eq:impl-dense-memory}; retaining
both construction and execution workspaces adds their sum, of the same
order. There is no hidden training phase, source-specific compilation,
or external arithmetic oracle.

\paragraph{A separate all-Boolean/ReLU unrolling.}
To remove the finite-attention primitive, select one complete
representation and unroll its fixed public bit-operation schedule. Let
$T$ be that representation's conservative online bitwork and $S$ its total
live bit storage, including input, program/constants, masks, and scratch.
Every conditional data operation is already a Boolean mux; loop lengths
and permitted routes are public. Replace each elementary bit operation by
\eqref{eq:impl-boolean}--\eqref{eq:impl-mux}, and at every instruction
explicitly hold every other live rail. Induction on instructions proves
that the resulting network state equals the original finite evaluator's
state. The final dyadic decode is also retained and charged.

This produces a \emph{PURE} Boolean/ReLU network, with no finite-attention
arithmetic primitive, having the valid bounds
\begin{align}
 \text{depth}&=O(T+p),&
 \text{width}&=O(S+p),\notag\\
 \text{sparse gates and holds}&=O((T+p)(S+p)),&&\label{eq:impl-pure}\\
 \text{dense parameter bits}&=O\bigl(p(T+p)(S+p)^2\bigr),&&\notag\\
 \text{dense bitwork}&=O\bigl(p^2(T+p)(S+p)^2\bigr).&&\notag
\end{align}
For example, one may use
$T=O(np^4)$ and $S=O(p^4+n\ell_a+np^2)$ from the sparse hybrid.
The sparse-record description can conservatively be written with
$O(p(T+p)(S+p))$ bits: all relevant indices have $O(p)$ bits, since
$\log n\le p$ and $\log p\le p$.
Writing and evaluating the PURE representation incurs its own displayed
storage and work costs. These large unrolling bounds do not inherit the
hybrid depth $O(p^2)$ or per-token width $O(p)$. The small hybrid shape
therefore must never be presented as the shape of the fully expanded
finite evaluator.

\subsection{Observation comparisons and the classical baseline}
\label{sec:impl-comparisons}

\paragraph{The count-only obstruction has an information scope.}
Theorem~\ref{thm:finite-window}, specialized to $r=2$, gives an optimal
budget-blind competitive factor between
\[
 2^{-30}\left(1+\log\frac{M+2}{2}\right)
 \quad\text{and}\quad
 2^{20}\left(1+\log\frac{M+2}{2}\right)
\]
for exact triple counts plus both endpoint pairs, at supplied $1/M$.
The lower tests the same map across budgets; it is not a logarithmic
penalty at every individual budget. Counts plus only the protected pair,
and any rounded versions of those observations, are weaker and inherit
the obstruction. The present half-position statistics change the
observation, so that comparison supplies neither an architecture-wide
lower bound nor a neural depth separation.

\paragraph{A one-pass counter on equal positional information.}
On the same attacked word, a classical program keeps four exact triples
of counters $(J_u,C_u,D_u)$ and a two-bit sliding window. When symbol
$Z_{i+2}$ arrives, the window identifies $u=(Z_i,Z_{i+1})$: increment
$J_u$, increment $C_u$ if $Z_{i+2}=1-Z_i$, and increment $D_u$ if
$i<\lfloor M/2\rfloor$. Then shift the two-bit window. After the last
symbol it contains the protected endpoint. Thus the program computes
exactly \eqref{eq:impl-counts} in one pass, without knowing $k$.

Its counters, loop index, and sliding window use $O(\log M)$ auxiliary
bits, with $O(M\log M)$ conservative bitwork for the scan. On the exact
counters, the same positive-denominator check, gate, and finite output
division give the risk in \eqref{eq:impl-risk} with $\eta=0$.
If output precision is charged, the division uses
$O((\log M+b+1)^2)$ active bitwork and
$O(\log M+b+1)$ live bits, with additional holds if fully layered.
Alternatively, round the exact counters once at the end using
\eqref{eq:impl-round}. This reproduces the attention bank coordinate for
coordinate, so it is also an equal-bank baseline, not merely an
asymptotically comparable estimator.

The input models must be matched. If the word is resident, charge the
same $M+2$ input bits to the counter and the attention program. If a
read-once input stream is allowed, state that permission explicitly; do
not compare the counter's auxiliary workspace with the attention
program's batch total memory. The counter has no attention mask-storage
cost. Consequently no neural statistical or computational advantage over
equal positional information is asserted. Sparse masks make the proposed
realization cheaper than generic dense attention evaluation, but do not
surpass this sequential baseline.

\paragraph{What changes relative to the earlier finite bank.}
The construction of \citet{BCJ001} already supplies eight full-triple and
four half-row normalized coordinates, public half-position flags,
uniform attention, endpoint acquisition, an isolated finite observation,
and a paid finite sequential evaluator with optional full unrolling. It
also already distinguishes fractional bank precision from arithmetic word
width and allows a final rounded probability. Its approximate-gate,
quantization, noise, and source-edit results establish the central
finite-bank precision phenomenon in that earlier interface.

The present implementation makes a different, explicit choice at several
points. Two singleton heads acquire the lagged bits and endpoint. The
local Boolean circuit emits four direct $J/C/D$ triples. A zero-score
head uses exact accumulation followed by one nearest rounding, rather
than the earlier evaluator's nested floors. The numeric fault is a
canonical replacement \emph{after} that rounding, rather than an additive
analog change before a final floor. Bit circuits, scalar-grid extraction,
restoring division, and dyadic output give the specified two-layer,
three-head hybrid with growing pointwise depth. Before rounding, the
full-triple and direct-row representations are linearly interconvertible;
after separately applied rounding they need not define the same
statistical experiment. Neither port theorem, especially a port-noise
lower bound, transfers merely because both banks have twelve labels.

Uniform indicator aggregation, positional copying, and finite arithmetic
compilation are standard ingredients
\citep{WeissGoldbergYahav2021,MerrillSabharwal2023,EdelmanEtAl2024,
RajaramanEtAl2024}. The construction here uses exactly accumulated
indicators followed by rounding of their aggregate, not rounding each
attention weight $1/M$ onto the $h$-grid. Results about rounding attention
probability vectors therefore concern a different numerical operation.
The explicit singleton structural masks likewise must not be confused
with dense-causal-only attention or finite-temperature approximations to
copying. No conclusion is asserted about training, learnability, minimal
head count, total-precision optimality, architecture-wide lower bounds,
pure-network depth advantages, empirical performance, or priority for
these standard mechanisms.
